% =====================================================================
%  c06.tex —— §7 A∞ 增广唯一的三角范畴之例（原书 p.54–56；PDF p.13 末–p.15 末）
%  编号公式：式 (9)（(H(C))^i(X,Y)=0，原书 p.55）、式 (10)（交换图表，原书 p.55）
%  定理类环境：定理 7.1、推论 7.2、推论 7.3
%  图表：1 幅（式 (10) 的交换图表，tikz-cd 重画）
%  说明：§7 标题在 PDF p.13 末（原书 p.54）；PDF p.16 首行已是 §8 标题，
%        故本块止于 PDF p.15 末（原书 p.56），不含 p.16 的任何内容。
% =====================================================================

% ===== 原书 p.54 =====
\section{关于\texorpdfstring{$\ainf$}{A∞}增广唯一的三角范畴之例}
\label{sec:c06}

对$\ainf$范畴$\cC$，$H(\cC)$构成分次范畴。把态射的合成视为$2$次$\ainf$结构，
即可把$H(\cC)$看作高阶（$3$次以上）$\ainf$结构为自明的极小$\ainf$范畴。
现在按这个意义，考虑与$H(\cC)$为$\ainf$拟同型的$\ainf$范畴$\cC$。

这类$\ainf$范畴源自有理同伦论（rational homotopy theory），把它称为形式的
（formal）是自然的。无论如何，对$\cC$应用极小模型定理（定理 2.3）可知：
可取到与$\cC$为$\ainf$拟同型的极小$\ainf$范畴$\cC$，它的高阶$\ainf$结构
$m'_{3},m'_{4},\dots$为自明。也就是说，$\cC$本质上没有高阶积。当$\cC$是
形式的$\ainf$范畴时，由引理 3.2 得三角范畴等价
$\Tr(\cC)\simeq\Tr(H(\cC))$，这意味着$\Tr(\cC)$只由$H(\cC)$决定。

\transnote{原文把与$\cC$为$\ainf$拟同型的极小模型仍记作$\cC$，与原本的$\cC$
同符号。按上下文，后一个$\cC$是极小模型；两者只是$\ainf$拟同型，并非同一个
范畴。}

% ===== 原书 p.55 =====
下面考察分次范畴$H(\cC)$上、以态射的合成为$2$次$\ainf$结构的极小$\ainf$范畴
结构有多少种。这样得到的极小$\ainf$范畴称为分次范畴$H(\cC)$的$\ainf$扩张
（extension）。把高阶$\ainf$积视为自明的$\ainf$范畴$H(\cC)$自身就是一个
$\ainf$扩张的例子，但一般还会有别的例子。若分次范畴$H(\cC)$的任意$\ainf$扩张
都与$\ainf$范畴$H(\cC)$为$\ainf$同型，则称分次范畴$H(\cC)$的$\ainf$扩张是
自明的（trivial）。（若分次范畴$H(\cC)$的$\ainf$扩张自明，则$\cC$必与$\ainf$
范畴$H(\cC)$为$\ainf$拟同型，即$\cC$是形式的。）

\transnote{「$\ainf$扩张自明」比「$\cC$是形式的」更强。后者只要求$\cC$的极小
模型的高次积为零；前者要求$H(\cC)$的任意$\ainf$扩张都与$H(\cC)$本身
$\ainf$同型。原文括号里只写了「自明 $\Longrightarrow$ 形式」这一方向，定理 7.1
的假设用的正是「自明」。}

以下考虑分次范畴$H(\cC)$的$\ainf$扩张自明的$\ainf$范畴$\cC$。例如考虑满足
$H(\cC)=H^{0}(\cC)$的$\ainf$范畴$\cC$，也就是说，对所有$X,Y\in\cC$，若
$i\neq0$，则
\begin{equation}
  (H(\cC))^{i}(X,Y)=0
  \label{eq:c06-9}
\end{equation}
成立。此时，由次数计算可知，$\cC$的极小模型的高阶$\ainf$积为自明，因此分次范畴
$H(\cC)$的$\ainf$扩张是自明的。

\transnote{式\eqref{eq:c06-9}的显示式省去了量词，完整条件是「对所有
$X,Y\in\cC$与所有$i\neq0$」，由前一句给出。若读成对一切$i$成立，就等于令
$H(\cC)$在$0$次也为零，条件失去内容。}

对分次范畴$H(\cC)$的$\ainf$扩张自明的$\ainf$范畴$\cC$，考虑三角范畴
$\cT:=\Tr(\cC)=\Tr(H(\cC))$的$\ainf$增广。假设存在与$\Tw(\cC)$不同的另一个
$\ainf$增广$\cD$，也就是说，假设存在三角范畴等价
$\phi:\cT=H^{0}(\Tw(\cC))\to H^{0}(\cD)$。这等价于存在（分次）范畴等价
$H(\Tw(\cC))\simeq H(\cD)$。取$\cD$的全子$\ainf$范畴$\cC'$，使这个等价的限制
给出全子范畴的等价$H(\cC)\simeq H(\cC')$，则存在$\ainf$等价序列
\[
  \cC\simeq H(\cC)\simeq H(\cC')\simeq\cC'.
\]
特别地，$H(\cC')\simeq\cC'$是$\ainf$拟同型。这由下述事实得出：分次范畴
$H(\cC')\simeq H(\cC)$的$\ainf$扩张自明，而以上同调为$H(\cC')$的$\ainf$范畴
$\cC'$是形式的。

另一方面，$\Tw(\cC')$是$\cD$的全子$\ainf$范畴，故记对应的$\ainf$函子为
$\iota:\Tw(\cC')\to\cD$。此时得到图表
\begin{equation}
\begin{tikzcd}[column sep=1.2em, row sep=1.7em]
  {} & \Tw(\cC) \arrow[r,"\simeq"] \arrow[d]
  & \Tw(\cC') \arrow[r,"\iota"] \arrow[d]
  & \cD \arrow[d] \\
  \cT \arrow[r,equal] & \Tr(\cC) \arrow[r,"\simeq"]
  & \Tr(\cC') \arrow[r,"H^{0}(\iota)"]
  & H^{0}(\cD)
\end{tikzcd}
\label{eq:c06-10}
\end{equation}
其中上排箭头是$\ainf$函子，下排箭头是三角函子，且图表在问题 4.1 的意义下交换。
左侧的交换图表由对上述$\ainf$拟同型序列应用引理 3.2 得到；右侧的交换图表由
$\iota$的定义自然得到。

现在$H^{0}(\iota)$是（全忠实的）三角函子。把它与三角范畴等价函子
$\Tr(\cC)\to\Tr(\cC')$的合成（图 10 下排函子的合成）记作
$\phi':\Tr(\cC)\to H^{0}(\cD)$。则三角函子$\phi'$在$H^{0}(\cC)\subset\Tr(\cC)$
上的限制，与$\phi$在$H^{0}(\cC)$上的限制一致。特别地，$H^{0}(\cC)$是
$\cT=\Tr(\cC)$的生成系，故$\phi$与$\phi'$自然等价。
% ===== 原书 p.56 =====
由此结合图 10 可知：$\phi'$，也就是$\phi$，具有提升$\Tw(\cC)\to\cD$。

总结以上，得：

\begin{theorem}
  设$\cC$是分次范畴$H(\cC)$的$\ainf$扩张为自明的$\ainf$范畴，则三角范畴
  $\Tr(\cC)$的$\ainf$增广是唯一的。
  \label{thm:c06-7.1}
\end{theorem}

这一事实隐含在 [12] 的论证中，不过 [12] 讨论的主要是其中性质更好的情形：
$H(\cC)\subset\Tr(\cC)$是三角范畴$\Tr(\cC)$中的强例外生成系
（full strongly exceptional collection）。

首先，作为分次范畴$H(\cC)$的$\ainf$扩张自明的$\ainf$范畴$\cC$的例子，考虑
上面提到的满足$H(\cC)=H^{0}(\cC)$者。此时，$\Tr(\cC)$由代数
\[
  A:=\bigoplus_{X,Y\in\cC}(H(\cC))^{0}(X,Y)
\]
决定。

另一方面，设$\cC$为一般的$\ainf$范畴（假设$H(\cC)$的各态射空间都是有限维的）。
当$\cC$的对象只有有限个、记为$E_{1},\dots,E_{n}\in\cC$时，若对任意$i$有
$H(\cC)(E_{i},E_{i})\simeq K$，且对任意$i>j$有
\[
  H(\cC)(E_{i},E_{j})=0,
\]
则$(E_{1},\dots,E_{n})$确定了三角范畴$\Tr(\cC)$的一个例外生成系（这几乎就是
定义本身 [1, 2], cf.\ [11]）。特别地，若例外生成系$(E_{1},\dots,E_{n})$还满足：
对任意$i<j$与$k\neq0$有
\[
  (H(\cC))^{k}(E_{i},E_{j})=0,
\]
则$(E_{1},\dots,E_{n})$是$\Tr(\cC)$中的强例外生成系（[1, 2], cf.\ [11]）。

\transnote{两类生成系的指标方向不同：定义例外生成系时，要求$i>j$的
$H(\cC)(E_{i},E_{j})$为零；条件加强为强例外时，还要求$i<j$且$k\neq0$的
$(H(\cC))^{k}(E_{i},E_{j})$为零。两处$i>j$与$i<j$的方向不要抄反。}

当这个$\cC$构成强例外生成系时，$\cC$是满足上述$H(\cC)=H^{0}(\cC)$者中性质相当
好的一类。此时$\Tr(\cC)$等价于有限维代数$A$上有限生成左模构成的导出范畴
$D^{b}(\mathrm{mod}\text{-}A)$（[1, 2], cf.\ [11]）。也就是说，此时可以断言：
导出范畴$D^{b}(\mathrm{mod}\text{-}A)$的$\ainf$增广是唯一的（[12, 推论 5.6]）。
在更一般的情形，若$A$是同调维数有限的有限维代数，同样有
$\Tr(\cC)\simeq D^{b}(\mathrm{mod}\text{-}A)$，这时的$\ainf$增广也是唯一的。

最后，列出两个由定理 7.1 容易得到的事实。

\begin{corollary}
  对满足$H(\cC)=H^{0}(\cC)$的$\ainf$范畴$\cC$，三角范畴$\Tr(\cC)$的
  $\ainf$增广是唯一的。
  \label{cor:c06-7.2}
\end{corollary}

\begin{corollary}
  当$\cC$是由$3$个以下对象构成的$\ainf$范畴、且构成$\Tr(\cC)$的例外生成系时，
  三角范畴$\Tr(\cC)$的$\ainf$增广是唯一的。
  \label{cor:c06-7.3}
\end{corollary}

事实上，分次范畴$H(\cC)$的$\ainf$扩张唯一：推论 7.2 的情形如上所述；推论 7.3
的情形也由不存在$3$个以上态射的$\ainf$积而显然。
